% Part: computability
% Chapter: recursive-functions
% Section: pr-functions

\documentclass[../../../include/open-logic-section]{subfiles}

\begin{document}

\olfileid{cmp}{rec}{prf}
\olsection{ఆదిమ పునరావృత్త ప్రమేయాలు}
\sourcecorrection{OLTECMPREC-003}{ఆంగ్ల మూల శీర్షికలో “Primitive Recursion Functions” అని ఉంది; అధ్యాయం నిర్వచించే ప్రమేయాల ప్రామాణిక పేరుకు అనుగుణంగా “Primitive Recursive Functions” అనే అర్థంతో అనువదించాం.}

ఆదిమ పునరావృత్తి, సంయుక్తం ఉపయోగించి ఇప్పటికే ఉన్న ప్రమేయాల నుంచి
కొత్త ప్రమేయాలను ఎలా నిర్వచించవచ్చో మళ్లీ నమోదు చేద్దాం.

\begin{defn}
  \ollabel{defn:primitive-recursion} $f$ ఒక $k$-స్థానిక ప్రమేయం
  ($k\ge 1$), $g$ ఒక $(k+2)$-స్థానిక ప్రమేయం అనుకుందాం. $f$, $g$ల
  నుంచి \emph{ఆదిమ పునరావృత్తి ద్వారా నిర్వచించిన ప్రమేయం} అంటే కింది
  సమీకరణాలతో నిర్వచితమైన $(k+1)$-స్థానిక ప్రమేయం~$h$:
  \begin{align*}
    h(x_0,\dots,x_{k-1},0) & =  f(x_0,\dots,x_{k-1}) \\
    h(x_0,\dots,x_{k-1},y+1) & = g(x_0,\dots,x_{k-1}, y, h(x_0,\dots,x_{k-1}, y))
  \end{align*}
\end{defn}

\begin{defn}
  \ollabel{defn:composition} $f$ ఒక $k$-స్థానిక ప్రమేయం; $g_0$,
  \dots, $g_{k-1}$ అన్నీ $n$-స్థానికాలైన $k$ ప్రమేయాలు అనుకుందాం.
  $f$, $g_0$, \dots,~$g_{k-1}$ల నుంచి \emph{సంయుక్తం ద్వారా
  నిర్వచించిన ప్రమేయం} అంటే కింది విధంగా నిర్వచితమైన $n$-స్థానిక
  ప్రమేయం~$h$:
  \[
  h(x_0, \dots, x_{n-1}) =
  f(g_0(x_0, \dots, x_{n-1}), \dots, g_{k-1}(x_0, \dots, x_{n-1})).
  \]
\end{defn}

$\Succ$, ప్రక్షేప ప్రమేయాలు
\[
\Proj{n}{i}(x_0,\dots,x_{n-1}) = x_i,
\]
తో పాటు, ప్రతి సహజ సంఖ్య $n$, ప్రతి $i<n$కూ, ప్రమేయం
~$\Zero(x)=0$ను ఆదిమ పునరావృత్త ప్రమేయాల్లో చేర్చుతాం.

\begin{defn}
  ఆదిమ పునరావృత్త ప్రమేయాల సమితి అంటే $\Nat^n$ నుంచి $\Nat$కు
  వెళ్లే ప్రమేయాల్లో, కింది నిబంధనలతో ఆగమనాత్మకంగా నిర్వచించిన సమితి:
  \begin{enumerate}
  \item $\Zero$ ఆదిమ పునరావృత్త ప్రమేయం.
  \item $\Succ$ ఆదిమ పునరావృత్త ప్రమేయం.
  \item ప్రతి ప్రక్షేప ప్రమేయం $\Proj{n}{i}$ ఆదిమ పునరావృత్త
    ప్రమేయం.
  \item $f$ ఒక $k$-స్థానిక ఆదిమ పునరావృత్త ప్రమేయం, $g_0$,
    \dots,~$g_{k-1}$లు $n$-స్థానిక ఆదిమ పునరావృత్త ప్రమేయాలైతే,
    $f$, $g_0$, \dots,~$g_{k-1}$ల సంయుక్తం ఆదిమ
    పునరావృత్త ప్రమేయం.
  \item $f$ ఒక $k$-స్థానిక ఆదిమ పునరావృత్త ప్రమేయం, $g$ ఒక
    $k+2$-స్థానిక ఆదిమ పునరావృత్త ప్రమేయమైతే, $f$, $g$ల నుంచి ఆదిమ
    పునరావృత్తి ద్వారా నిర్వచించిన ప్రమేయం ఆదిమ పునరావృత్త ప్రమేయం.
\end{enumerate}
\end{defn}

\begin{explain}
ఇంకా క్లుప్తంగా చెప్పాలంటే, ఆదిమ పునరావృత్త ప్రమేయాల సమితి
$\Zero$, $\Succ$, ప్రక్షేప ప్రమేయాలు~$\Proj{n}{j}$లను కలిగి,
సంయుక్తం, ఆదిమ పునరావృత్తి కింద సంవృతంగా ఉండే అతి చిన్న సమితి.

ఆదిమ పునరావృత్త ప్రమేయాల సమితిని ``దశల'' పరంగా వర్ణించడం మరొక
పద్ధతి. ప్రారంభ ప్రమేయాలైన $\Zero$, $\Succ$, ప్రక్షేపాల సమితిని
$S_0$ అనుకుందాం. ఇవి దశ~$0$ ఆదిమ పునరావృత్త ప్రమేయాలు. ఒకసారి దశ
$S_i$ నిర్వచితమైతే, $S_{i+1}$లో $S_i$లోని ప్రమేయాలన్నిటినీ, అలాగే
$S_i$లో ఇప్పటికే ఉన్న ప్రమేయాలకు సంయుక్తం లేదా ఆదిమ పునరావృత్తి యొక్క
ఒక్క ప్రయోగంతో లభించే ప్రమేయాలన్నిటినీ చేర్చుదాం. అప్పుడు
\[
S = \bigcup_{i \in \Nat} S_i
\]
అన్ని ఆదిమ పునరావృత్త ప్రమేయాల సమితి.
\sourcecorrection{OLTECMPREC-004}{ఆంగ్ల మూలం $S_{i+1}$లో కేవలం $S_i$ ప్రమేయాలకు ఒక్క నిర్మాణ చర్యతో లభించిన ప్రమేయాలనే ఉంచి, $S_i$నే చేర్చలేదు; భిన్న నిర్మాణ లోతుల ప్రమేయాలను తరువాత కలపడానికి అవసరమైన సంచిత దశ నిర్వచనాన్ని స్పష్టంగా చేర్చాం.}
\end{explain}

$\Add$ ఒక ఆదిమ పునరావృత్త ప్రమేయమని నిర్ధారిద్దాం.

\begin{prop}
  కూడిక ప్రమేయం $\Add(x,y)=x+y$ ఆదిమ పునరావృత్త ప్రమేయం.
\end{prop}

\begin{proof}
$\Add$కు $f$,~$g$ అనే రెండు ప్రమేయాల పరంగా ఇప్పటికే ఉన్న ఆదిమ
పునరావృత్త నిర్వచనం \olref{defn:primitive-recursion} రూపానికి
సరిపోతుంది:
\begin{align*}
  \Add(x_0, 0) & = f(x_0) = x_0\\
  \Add(x_0, y+1) & = g(x_0, y, \Add(x_0, y)) =
  \Succ(\Add(x_0, y))
\end{align*}
కాబట్టి $f$, $g$లు ఆదిమ పునరావృత్త ప్రమేయాలైతే $\Add$ కూడా అలాంటిదే.
$f(x_0)=x_0=\Proj{1}{0}(x_0)$; ప్రక్షేప ప్రమేయాలు ఆదిమ పునరావృత్త
ప్రమేయాలే కనుక $f$ కూడా అలాంటిదే. $g$ అనేది $g(x_0,y,z)$ అనే
\[
g(x_0, y, z) = \Succ(z).
\]
అని నిర్వచించిన త్రిస్థానిక ప్రమేయం. దీనివల్ల మాత్రమే $g$ ఆదిమ
పునరావృత్త ప్రమేయమని ఇంకా తేలదు; ఎందుకంటే $g$, $\Succ$ సరిగ్గా ఒకే
ప్రమేయం కావు: $\Succ$ ఏకస్థానికం, $g$ త్రిస్థానికం కావాలి. అయితే
సంయుక్తం ద్వారా $g$ను నియమబద్ధంగా
\[
g(x_0, y, z) = \Succ(\Proj{3}{2}(x_0, y, z))
\]
అని నిర్వచించవచ్చు. $\Succ$, $\Proj{3}{2}$ ఆదిమ పునరావృత్త
ప్రమేయాలు కనుక, వాటి సంయుక్తంతో నిర్వచితమైన $g$ కూడా ఆదిమ పునరావృత్త
ప్రమేయమే.
\end{proof}

\begin{prop}
  \ollabel{prop:mult-pr}
  గుణకార ప్రమేయం $\Mult(x,y)=x\cdot y$ ఆదిమ పునరావృత్త ప్రమేయం.
\end{prop}

\begin{proof}
  అభ్యాసం.
\end{proof}

\begin{prob}
  \olref[cmp][rec][prf]{prop:mult-pr}ను నిరూపించడానికి, $\Mult$ యొక్క
  ఆదిమ పునరావృత్త నిర్వచనాన్ని
  \olref[cmp][rec][prf]{defn:primitive-recursion} కోరే రూపంలో
  రాయవచ్చని, దానికి సంబంధించిన $f$,~$g$ ప్రమేయాలు ఆదిమ పునరావృత్త
  ప్రమేయాలని చూపండి.
\end{prob}

\begin{ex}
ఆదిమ పునరావృత్త నిర్వచనానికి మన మొట్టమొదటి ఉదాహరణ ఇదే:
\begin{align*}
h(0) & =  1 \\
h(y+1) & =  2 \cdot h(y).
\end{align*}
$k=0$ కనుక ఈ ప్రమేయం \olref{defn:primitive-recursion} కోరే రూపంలో
ఇమడదు. నిర్వచనంలో స్థిరాంకాలు $1$, $2$ కూడా ఉన్నాయి. మొదటి సమస్యను
దాటడానికి, ఒక పూరక ఆర్గుమెంటును ప్రవేశపెట్టి ప్రమేయం~$h'$ను
నిర్వచిద్దాం:
\begin{align*}
h'(x_0, 0) & =  f(x_0) = 1 \\
h'(x_0, y+1) & =  g(x_0, y, h'(x_0, y)) = 2 \cdot h'(x_0, y).
\end{align*}
$f(x_0)=1$ను $\Succ$, $\Zero$ల నుంచి సంయుక్తం ద్వారా నిర్వచించవచ్చు:
$f(x_0)=\Succ(\Zero(x_0))$. $g$ను $g'(z)=2\cdot z$, ప్రక్షేపాల
నుంచి సంయుక్తం ద్వారా నిర్వచించవచ్చు:
\begin{align*}
g(x_0, y, z) & = g'(\Proj{3}{2}(x_0, y, z))
\intertext{$g'$ను కూడా సంయుక్తం ద్వారా కింది విధంగా నిర్వచించవచ్చు:}
g'(z) & = \Mult(g''(z), \Proj{1}{0}(z))
\intertext{మరియు}
g''(z) & = \Succ(f(z)),
\end{align*}
ఇక్కడ పైన ఉన్న $f$కు $f(z)=\Succ(\Zero(z))$. ఇప్పుడు $h'$ మనకు ఉంది
కనుక, మళ్లీ సంయుక్తాన్ని ఉపయోగించి $h(y)=
h'(\Proj{1}{0}(y),\Proj{1}{0}(y))$గా తీసుకోవచ్చు. ఈ విధంగా ప్రాథమిక
ప్రమేయాల నుంచి సంయుక్త నిర్మాణాలు, ఆదిమ పునరావృత్తుల క్రమాన్ని ఉపయోగించి $h$ను
నిర్వచించవచ్చు; కనుక $h$ ఆదిమ పునరావృత్త ప్రమేయం.
\end{ex}

\end{document}
